\documentclass[12pt]{report}
\usepackage[a4paper, margin=2.5cm]{geometry}
\pagenumbering{gobble}
\usepackage{hyperref}
\usepackage{tcolorbox}
\newcommand{\student}[3]{\begin{tcolorbox}\noindent Introduction to
    Psycholinguistics (WS2021-2022) \\ Report of \textbf{Paper~#3} by
    \textbf{#1 (#2)} 
  \end{tcolorbox}}



\begin{document}

%% Change:
\student{Shawon Ashraf}{3460184}{3}


\paragraph{APA reference:} 
%report the full reference of the paper in APA style (\href{https://www.library.kent.edu/files/APACheatSheet.pdf}{\underline{Online Guide}})
Trueswell, J. C. (1996). The role of lexical frequency in syntactic ambiguity resolution. \textit{Journal of memory and language}, 35(4), 566-585.

\paragraph{Research Question(s):}
%What are the questions the authors of the paper are addressing?

\begin{enumerate}
    \item Does the frequency of a verb being used as a participle form influence a reader to consider the alternative of a relative clause?
    
    \item Do readers disambiguate syntactic and lexical ambiguities in the same fashion? Or, continuing from the last question, does lexical frequency affect syntactic ambiguity resolution?
    
    \item How much effect does relative frequency of a relative clause has over to its overall frequency when compared to main clauses?
\end{enumerate}

\vfill
\paragraph{Methodology:} 
%which methods are used?

The author divided their experiment in two parts. \textbf{In the first part}, they used moving-window self-paced reading to check how lexical frequency can affect syntactic ambiguity resolution. Two different groups of verbs were used in a self paced read one word at a time task. Half of the verbs had high participle frequency and were called High-PP Verbs while the other half were called Low-PP Verbs for having the opposite amount of participle frequency. All the verbs used in the experiment were morphologically ambiguous. The verbs were paired with nouns to set up an Agent - Patient theme. 

There were fourteen subjects in the experiment, all from University of Pennsylvania. The duration of the first experiment was set to 30 minutes. The subjects in the experiment were also asked to complete sentences based on these verbs to test if a verb's participle frequency can affect the choice of a clause during sentence generation. Half of the subjects attended this study. 

\textbf{The second experiment}  was conducted with the same materials from the first experiment. However, this time in noun-verb pairing they only used nouns which were good agents. The subjects were then asked to rate the sentences based on how good the agent-patient assignment was in a sentence. 

Eighteen subjects from University of Pennsylvania attended the experiment as subjects.
\vfill

\paragraph{Results:} 
%what are the main results?

\vfill

\paragraph{Personal Opinion:} 
%What did I like, dislike about this
%paper? 
\vfill


\end{document}

